\documentclass[12pt]{article}
\usepackage[utf8]{inputenc}
\usepackage[T1]{fontenc}
\usepackage{hyperref}
\title{Yoloxóchitl Mixtec Language Resources for K-12 Curriculum Development}
\author{Compiled by Claude Sonnet 5.5 (claude-sonnet-5-5) for Andr\'as Kornai}
\date{}
\begin{document}
\maketitle

\section{Introduction}
Yoloxóchitl Mixtec (ISO 639-3 xty; endonym reported as Tu'un savi) is an Otomanguean language of Guerrero, Mexico. It is not known to be taught in schools. This survey is based on a small number of web searches. The resources found are mainly documentation materials, and no school materials were found.

\section{Primary Educational Resources}
No primers, textbooks or school curricula for Yoloxóchitl Mixtec turned up in the searches. The documentation corpora listed below may be useful as source material for teachers and curriculum writers.

\section{Linguistic and Reference Materials}
\begin{itemize}
\item Open Speech and Language Resources, Yoloxóchitl Mixtec speech and transcription corpus, compiled by Jonathan Amith and Rey Castillo García. The listing gives the license as CC BY-NC-SA 3.0 and the archive size as about 86 GB. \url{https://www.openslr.org/89/}
\item AILLA (UT Austin), Yoloxóchitl Mixtec Language Documentation Project of Jonathan Amith and Rey Castillo. \url{https://islandora-ailla.lib.utexas.edu/islandora/object/ailla%3A124484}
\item Endangered Languages Archive, Yoloxochitl (Guerrero) Mixtec stories and other oral traditions. \url{https://eldp.access.preservica.com/dk0078/}
\item Oto-Manguean languages at Surrey, language information page for xty. \url{https://oto-manguean.surrey.ac.uk/Info/XTY}
\item Wikipedia, Yoloxóchitl Mixtec (overview). \url{https://en.wikipedia.org/wiki/Yolox%C3%B3chitl_Mixtec}
\end{itemize}
No dedicated dictionary or grammar appeared in the search results.

\section{Religious and Community Resources}
\begin{itemize}
\item Global Recordings Network, Mixtec, Yoloxochitl: audio Bible-based and evangelistic recordings. \url{https://globalrecordings.net/en/language/xty}
\end{itemize}

\section{Digital and Audio Resources}
\begin{itemize}
\item Wikitongues, Yoloxóchitl Mixtec language page. \url{https://wikitongues.org/languages/xty/}
\item The OpenSLR, AILLA and ELAR collections above contain audio recordings with transcriptions.
\end{itemize}

\section{Curriculum Development Context}
The language is reported as stable and in use by the whole community, but it is not known to be taught in schools. The audio and transcribed documentation by Amith and Castillo is the most substantial resource found. It would need community and researcher involvement, and attention to the license terms, before it could be adapted for classroom use. A K-12 curriculum would have to be built almost from scratch.
\end{document}